\documentclass[11pt,a4paper,oneside,openany]{book}

\usepackage[utf8]{inputenc}
\usepackage[T1]{fontenc}
\usepackage[english]{babel}
\usepackage{csquotes}
\usepackage[charter]{mathdesign}
\usepackage{amsmath}
\usepackage{lipsum}
\usepackage{amsthm}
\usepackage{mathtools}
\usepackage{physics}
\usepackage{graphicx}
\usepackage{enumitem}
\usepackage{tikz-cd}
\usepackage{geometry}

% Page header management
\usepackage{fancyhdr}
\pagestyle{fancy}
\fancyhf{} 
\fancyhead[L]{\slshape\nouppercase{\leftmark}} 
\fancyhead[R]{\thepage} 
\renewcommand{\headrulewidth}{0pt}

% Redefinition of the plain style for chapter starting pages
\fancypagestyle{plain}{
	\fancyhf{}
	\renewcommand{\headrulewidth}{0pt}
	\fancyfoot[R]{\thepage} 
}

\usepackage[backend=biber, sorting=none, style=alphabetic]{biblatex}
\addbibresource{bib.bib}
\usepackage[colorlinks=true,linkcolor=black,citecolor=black,urlcolor=black]{hyperref}

\usepackage{epigraph}
\setlength{\epigraphwidth}{0.5\textwidth}
\renewcommand{\epigraphflush}{flushright}
\renewcommand{\textflush}{flushright}

\newtheorem{theorem}{Theorem}[chapter]
\newtheorem{definition}{Definition}[section]
\newtheorem{lemma}[definition]{Lemma}
\newtheorem{proposition}[definition]{Proposition}
\newtheorem{corollary}[definition]{Corollary}
\theoremstyle{remark}
\newtheorem{remark}[definition]{Remark}
\newtheorem{example}[definition]{Example}

\let\mathcal\relax
\newcommand{\mathcal}[1]{\text{\usefont{OMS}{cmsy}{m}{n}#1}}
\DeclareMathOperator{\Spec}{Spec}
\DeclareMathOperator{\Ad}{Ad}
\DeclareMathOperator{\ad}{ad}
\newcommand{\poisson}[2]{\{#1,#2\}}
% ============================================================
% Comandi Matematici
% ============================================================

% --- Insiemi numerici ---
\newcommand{\R}{\mathbb{R}}
\newcommand{\C}{\mathbb{C}}
\newcommand{\N}{\mathbb{N}}
\newcommand{\K}{\mathbb{K}}
\newcommand{\A}{\mathfrak{A}}
\newcommand{\B}{\mathcal{B}}

% --- Spazio di Hilbert e spazio proiettivo ---
% ATTENZIONE: \H è già un comando di LaTeX (accento ungherese, es. \H{o}=ő).
% Il renewcommand lo sovrascrive: nella tesi non ti servirà mai quell'accento,
% quindi è sicuro, ma se preferisci non toccare comandi di sistema usa \Hil
% al posto di \H ovunque sotto.
\renewcommand{\H}{\mathcal{H}}
\newcommand{\PH}{\mathbb{P}\H}                    % P H, spazio proiettivo

% --- Algebre C* ---
\newcommand{\Cstar}{C^*}                          % da usare in math mode: $\Cstar$
\newcommand{\Csalg}{$C^*$-algebra}                % da usare nel testo: "è una \Csalg\ commutativa"
\newcommand{\Csalgs}{$C^*$-algebre}               % plurale
\newcommand{\BH}{\mathcal{B}(\H)}                           % operatori limitati su H
\newcommand{\Cz}[1]{C_0(#1)}                      % C_0(M)
\newcommand{\Cinf}[1]{C^\infty(#1)}               % C^infty(M)
\newcommand{\Cinfc}[1]{C^\infty_c(#1)}            % C^infty_c(M)
\newcommand{\Mat}[2]{M_{#1}(#2)}                  % M_n(C), es. \Mat{2}{\C}

% --- Algebre di Lie (fraktur generico) ---
\newcommand{\mf}[1]{\mathfrak{#1}}                % \mf{g}, \mf{su}(2), \mf{u}(n), \mf{X}(M)

\newcommand{\qcomm}[2]{\dfrac{1}{i\hbar}[#1,#2]}  % (1/ih)[A,B]  -- condizione di Dirac, Qubit IV

% --- Notazione bra-ket ---
\newcommand{\melem}[3]{\langle#1|#2|#3\rangle}    % elemento di matrice <psi|A|psi>
\newcommand{\expect}[2]{\langle#1\rangle_{#2}}    % <A>_psi

% --- Campo vettoriale hamiltoniano ---
\newcommand{\Ham}[1]{X_{#1}}                      % \Ham{f}, \Ham{f_H}, \Ham{f_A}

% --- Simboli e operatori vari ---
\newcommand{\Id}{\mathbb{1}}
\renewcommand{\Re}{\mathrm{Re}}                   % la tesi usa Re/Im in tondo, non ℜ/ℑ
\renewcommand{\Im}{\mathrm{Im}}
\newcommand{\LC}{\varepsilon_{ijk}}               % simbolo di Levi-Civita (indici fissi ijk)
\DeclareMathOperator{\Ran}{Ran}
\DeclareMathOperator{\Ker}{Ker}
\DeclareMathOperator{\Span}{span}
\DeclareMathOperator{\sing}{sing}
\newgeometry{margin=2cm}
\setlength{\headwidth}{\textwidth}
\begin{document}
\setcounter{chapter}{2}
\epigraph{Nobody takes seriously quantum mechanics.}{Adán Cabello Quintero}

Quantum phenomena force a redefinition of physical reality and thus require a new descriptive language. In fact, words like 'to measure' are used but their operational semantics shifts. This mirrors the adaptation of Newtonian terms like time, mass, and energy in special relativity. At the same time, some syntactically correct propositions lose physical meaning. An example in special relativity is the assertion that two events occurred simultaneously at different places.

In order to find this new framework, we start by recalling that a physical theory partitions the world into the system and its environment. The framework described here treats only the system while the environment remains classical. This duality is underlined by Niels Bohr \cite{schilppAlbertEinsteinPhilosopherscientist1991}.
\begin{quote}
	\textit{However far the [quantum] phenomena transcend the scope of classical physical explanation, the account of all evidence must be expressed in classical terms. The argument is simply that by the word 'experiment' we refer to a situation where we can tell others what we have done and what we have learned and that, therefore, the account of the experimental arrangement and the results of the observations must be expressed in unambiguous language with suitable application of the terminology of classical physics.}
\end{quote}
Bohr used to reinforce this point with pseudorealistic sketches of robust clocks and rigid supports. It must be highlighted that he did not claim different physical laws for microscopic and macroscopic systems. He insisted on distinct modes of description between system and measurement apparatus. This approach is nonetheless \emph{unsatisfactory}. Motivated by these considerations, the main aim of the present work is to show the common algebraic structure underlying these two theories understanding the common grounds and the differences.

\section{General Mathematical Description of a Quantum Theory}\label{sec:intro-alg-quant}
This section shows that the same algebraic structure seen in the Hamiltonian case is forced on us directly by the operational meaning of measurement, without assuming a phase space
and without assuming that the resulting algebra is commutative. A system is specified solely by the states in which it can be prepared and the observables that can be measured on it. Starting from this minimal data, we build the observable set together with the norm and algebraic structure that let the observable set be embedded in an associative $C^*$-algebra. We will see that commutativity not used in the classical case is needed to encode the features of quantum mechanics.

We seek this new mode of description by analysing the procedures used in experimental physics. In fact, two categories arise, preparations and tests. In this approach, these are primitive undefined notions of quantum theory, analogous to points and lines in Euclidean geometry. In particular, a preparation is a completely specified and unambiguous experimental procedure. A measurement is an apparatus interacting with the system where a system property affects an apparatus property, and generally, this can disturb other properties of the system. That occurs even in classical physics, as shown in \cite[Section~12-2]{peresQuantumTheoryConcepts2002}. However, in that framework we assume the measured property exists prior to the interaction. Instead quantum physics rejects this premise, as shown in \cite[Section~1.5]{peresQuantumTheoryConcepts2002}. However, this epistemological discussion lies outside our present scope. We refer the reader to \cite{jammerPhilosophyQuantumMechanics1974} for a philosophical discussion about this topic.

Having set preparations and tests, we can now turn to define a 'quantum system'. Following N. Bohr \cite{bohrCanQuantumMechanicalDescription1935}, a complete description of a quantum system necessarily includes the measurements to be performed, as these determine the relevant physical quantities. A quantum system thus comprises a physical system, a set of preparations, and a set of measurements. The equivalence of distinct macroscopic procedures giving rise to the same physical object, such as a polarized photon or a free hydrogen atom, induces the notion of a quantum state, characterized by the probability distribution of measurement outcomes arising from repeated identical preparations. Likewise, an observable is defined as an equivalence class of measurements, where two measurements are identified if they yield identical outcome statistics across all the given preparations. This formal framework is developed further in the subsequent sections, where the property that an equivalence class uniquely determines the physical quantity is formalized as Separation in Definition \ref{def:state-space}.

\subsection*{The Observable Set and States}

An observable is operationally tied to a concrete measuring apparatus,
namely a device that, applied to any preparation of the system, returns a
real number. Two elementary manipulations of an apparatus carry a physical meaning. Rescaling the readout by a real factor $\lambda$, such as relabelling the pointer scale, produces a new
observable $\lambda A$ from $A$. Replacing every reading by its $n$-th
power produces a new observable $A^n$, obtained by squaring the scale,
cubing it, and so on, and this operation inherits the elementary law of
exponents, $A^mA^n=A^{m+n}$, directly from the arithmetic of real
numbers. In particular, $A^0$ is the trivial observable whose outcome is
always $1$, independently of the preparation. A quantum system is
specified jointly by a set $S$ of states and
by the set of observables obtainable in this way, which we now formalize.

\begin{definition}[Observable set]\label{def:obs-set}
	We call $\mathcal{O}$ an observable set if it is a set equipped
	with the operations
	\[
	(\lambda,A)\longmapsto \lambda A, \qquad (n,A)\longmapsto A^n,
	\qquad A^mA^n=A^{m+n}, \qquad A^0\equiv 1,
	\]
	for $\lambda\in\mathbb{R}$, $m,n\in\mathbb{N}$, $A\in\mathcal{O}$.
\end{definition}

An observable is called \emph{positive} if every possible outcome
of its measurement is a nonnegative number. Operationally, this is realized when the observable is obtained by squaring another one, for instance by taking a squared pointer scale, whose readings are nonnegative.

\begin{definition}[Positive observable]
	\label{def:positive-obs}
	Let $A\in\mathcal{O}$ with $\mathcal{O}$ denoting an observable set as in Definition \ref{def:obs-set}, $A$ is \emph{positive} if $A=B^2$ for some
	$B\in\mathcal{O}$.
\end{definition}

Consider the set of preparations given with a quantum system. Through the procedure described earlier, the equivalence classes of identical preparations are called \emph{states}, denoted by $\omega$. The expectation value $\omega(A) \in \mathbb{R}$ is obtained as the limit of sample means of replicated measurements of $A$ on identically prepared copies of the system. Since sample means are linear in the individual readings and intrinsically non-negative when evaluating squares of outcomes, we elevate these operational facts to defining axioms for the states.

	\begin{definition}[State]\label{def:q-state}
		A \emph{state} $\omega$ on the observable set $\mathcal{O}$, as per Definition \ref{def:obs-set}, is a real linear functional $\omega: \mathcal{O} \to \R$ satisfying two properties:
		
		\begin{itemize}
			\item Normalization: $\omega(\Id) = 1$.
			\item Positivity: $\omega(B^2) \ge 0$ for every $B \in \mathcal{O}$.
		\end{itemize}
	\end{definition}
	
	\noindent Having defined states, we now present the state space $S$, which, in complete analogy with the classical case, is convex.
	
	\begin{definition}[State space]\label{def:state-space}
		The \emph{state space} $S$ of a system is a convex set of states which is assumed to be \emph{separating} for $\mathcal{O}$. To wit, for any two observables $A,B\in\mathcal{O}$, if $\omega(A)=\omega(B)$ for all $\omega\in S$, then $A=B$.
	\end{definition}
	
	The separating property mathematically implements the operational completeness of the state space. The physical content of an observable is entirely exhausted by its expectations across all possible preparations. Two states that agree on every observable are identified as one and the same state. This mutual duality is  symmetric.
	\begin{remark}
		Definition \ref{def:q-state} guarantees by construction that every state evaluates non-negatively on the algebraic squares of observables. However, the converse implication, namely that state-positivity guarantees the existence of an exact square root within $\mathcal{O}$, fails. Establishing this strict equivalence requires endowing the observable set with a topological structure, such as norm completeness within a $C^*$-algebra. In this setting the continuous functional calculus ensures that such roots are always well defined.
	\end{remark}

Every concrete apparatus has an inevitable scale bound, for instance a
bounded pointer range, independent of the state being probed. Consequently,
the outcomes of measuring a fixed observable $A$, across all possible
states of Definition~\ref{def:q-state}, form a bounded set of real numbers,
and it is natural to record this bound as a norm.

\begin{definition}[Norm on $\mathcal{O}$]\label{def:normed-obs-set}
	Let $\mathcal{O}$  be the observable set defined in Definition \ref{def:obs-set} and let $S$ be the associated state space as in Definition \ref{def:state-space}. We can endow $\mathcal{O}$ with a norm $\| \cdot \| : \mathcal{O} \to \R$ defined as 
	\[
	\|A\| := \sup_{\omega \in S} |\omega(A)| < \infty .
	\]
\end{definition}

This norm is an additional structure layered on the bare set of
Definition~\ref{def:obs-set}. Note that we continue to write $\mathcal{O}$ for the
resulting normed set, rather than introduce new notation at each
enrichment, and will do so consistently from here on. 

By homogeneity of the states, $\|\lambda A\|=|\lambda|\,\|A\|$. Since the states separate
the observables, $\|A\|=0$ implies $A=0$. We can also prove that using only
positivity of the states applied to $\|A\| \Id \pm A$ and to
$(\|A\| \Id \pm A)^2$, the norm can be linked to the squaring operation already
available on $\mathcal{O}$.

\begin{proposition}[Compatibility of norm and squares]\label{prop:comp-squares}
	Let $\mathcal{O}$ be the observable set defined in Definition \ref{def:normed-obs-set}. For every $A\in\mathcal{O}$, $\|A\|^2=\|A^2\|$.
\end{proposition}

The complete proof of this statement can be found in \cite[p.18]{strocchiIntroductionMathematicalStructure2005}. Operationally, the sum of two observables is meaningful whenever there is
a third observable whose expectation always equals the sum of the two
expectations, for instance kinetic and potential energy. Postulating the
existence of such sums for arbitrary pairs turns $\mathcal{O}$ into a real
vector space, on which $\|\cdot\|$ is subadditive. Its completion is a real
Banach space, still denoted $\mathcal{O}$.

\begin{definition}[Linear extension of $\mathcal{O}$]
	Let $\mathcal{O}$ be an observable set as in Definition \ref{def:normed-obs-set} and let $S$ be the associated state space as in Definition \ref{def:state-space}. For any observables $A,B,C \in \mathcal{O}$, if the identity $\omega(C)=\omega(A)+\omega(B)$ holds for all states $\omega \in S$, we define the sum as $C=A+B$. 
	
	The linear extension of $\mathcal{O}$ is postulated by requiring that such sums, alongside their induced powers $(A+B)^n$, exist for all arbitrary pairs within $\mathcal{O}$. Furthermore the states in $S$ must act upon these extended elements as positive linear functionals.
\end{definition}

The structural consequence of this operational hypothesis is that $\mathcal{O}$ acquires the algebraic structure of a real vector space. Furthermore, this linear framework is compatible with the previously defined topological properties. Specifically, the norm becomes subadditive, guaranteeing that
$$\|A+B\| \le \|A\| + \|B\|, \qquad \|\omega(A)\| \le \|A\|,$$
for all $A, B \in \mathcal{O}$ and all $\omega \in S$. Consequently, $\mathcal{O}$ constitutes a normed real vector space. Taking the norm completion of this space yields a real Banach space, providing the necessary topological completeness for the subsequent algebraic development. This is a strictly larger structure that we continue to denote $\mathcal{O}$.

\subsection*{The Jordan product}

So far $\mathcal{O}$ carries only the Banach space structure just
introduced, together with the squaring of a single observable. To reach
an associative algebra, and eventually the full apparatus of operator theory, a product combining two different observables is needed. Nothing in the operational setting singles out a product $AB$ between two \emph{arbitrary} observables, since only squares, and hence sums of
squares, carry a direct experimental meaning. The symmetric
combination built from the operations already available, namely squares
and sums, is the following.

\begin{definition}[Jordan product]
	\label{def:jordan-product}
	Let $\mathcal{O}$ be an observable set, as in Definition \ref{def:obs-set}. The Jordan product is a commutative operation between two observables $A, B \in \mathcal{O}$ defined via
	\[
	A\circ B := \tfrac12\big((A+B)^2-A^2-B^2\big) = B\circ A .
	\]
\end{definition}

\begin{remark}
	As with the earlier enrichments, we do not introduce new notation for
	this product. For polynomials of a single observable $A$, the Jordan Product $\circ$
	coincides with the ordinary algebraic product, so that $A\circ
	A^n=A^{n+1}$.
\end{remark}

At this stage, $\circ$ is not known to be distributive over the sum
already available on $\mathcal{O}$. This property follows, however, from
a mild extra assumption, homogeneity of $\circ$, itself a direct extension
of the homogeneity already required on scalars.

\begin{proposition}[Homogeneity implies distributivity]
	Let $\mathcal{O}$ be an observable set with a Jordan product $\circ$ as in Definition \ref{def:jordan-product}. If $\circ$ is homogeneous, namely $A\circ(\lambda B)=\lambda(A\circ B)$ for
	$\lambda\in\mathbb{R}$, then $\circ$ is distributive,
	$A\circ(B+C)=A\circ B+A\circ C$.
\end{proposition}

This is proved in \cite[p.20]{strocchiIntroductionMathematicalStructure2005}. Distributivity alone says nothing about the norm of $A\circ B$. Only combining it with the positivity of the states and Proposition \ref{prop:comp-squares} enforces a bound on $\| A\circ B \|$, as shown in the following proposition.

\begin{proposition}[Norm bound on the Jordan product]
	Let $\mathcal{O}$ be an observable set with Jordan product $\circ$ as in Definition \ref{def:jordan-product}. Then $\|A\circ B\|\le\|A\|\,\|B\|$ for all $A,B\in\mathcal{O}$.
\end{proposition}
\begin{proof}
	For every state $\omega$ and every $\lambda\in\mathbb{R}$, positivity of
	$(A+\lambda B)^2$ yields
	$$0\le\omega(A^2)+2\lambda\,\omega(A\circ B)+\lambda^2\omega(B^2)\le
	\|A\|^2+2\lambda\,\omega(A\circ B)+\lambda^2\|B\|^2.$$ A nonnegative
	discriminant for this quadratic in $\lambda$ forces
	$\omega(A\circ B)^2\le\|A\|^2\|B\|^2$ for every $\omega$. This yields
	$\|A\circ B\|\le\|A\|\,\|B\|$.
\end{proof}

One algebraic property of a Jordan algebra remains open, it is called weak
associativity $$A^2\circ(B\circ A)=(A^2\circ B)\circ A.$$ This key identity ensures power associativity, guaranteeing that arbitrary polynomials of a single observable are unambiguously defined. Algebraically, this provides the main structure required to establish the continuous functional calculus. The commutativity, distributivity, weak associativity, and norm properties established above are, taken together, exactly the defining axioms of a
single algebraic structure. Unlike the earlier enrichments, we now give this structure its own name and record its axioms once and for all in the following definition.

\begin{definition}[Jordan--Banach algebra]
	\label{def:jb-algebra}
	A Jordan--Banach algebra $\mathcal{O}$ is a real Banach space as in Definition \ref{def:banach-space}, equipped with a commutative, distributive, weakly associative product $\circ$ such that
	\[
	\|A\circ B\|\le\|A\|\,\|B\|, \qquad \|A\|^2=\|A^2\|,
	\qquad \|A^2\|\le\|A^2+B^2\| \qquad \forall A,B \in \mathcal{O}.
	\]
	Such a structure is also called a \emph{JB algebra}.
\end{definition}

\subsection*{Embedding a JB Algebra into an Associative Algebra}

The Jordan product alone is, in general, neither associative nor known to
arise from an ordinary matrix-like product. It is technically convenient, and possible for every physically reasonable observable set, to realize
$\mathcal{O}$ concretely inside an associative $C^*$-algebra, so that
$\circ$ becomes the symmetrization of the usual product.

\begin{definition}[Embedding of a JB algebra in a $C^*$-algebra]\label{def:jb-embedding}
	A Jordan--Banach algebra $\mathcal{O}$, as in Definition~\ref{def:jb-algebra},
	is \emph{embedded} in a $C^*$-algebra $\A$ if there exists an isometric linear map from $\mathcal{O}$ into $\A$ such that the image of $\mathcal{O}$ generates $\A$ as a $C^*$-algebra. In addition, for all $A,B\in\mathcal{O}$
	\begin{enumerate}
		\item $A^* = A$ in $\A$, so that $\mathcal{O}$ is identified with a closed real subspace of the selfadjoint elements $\A_{sa}$ as per Definition \ref{def:selfadj-element},
		\item the Jordan product corresponds to the symmetrized associative product, $A\circ B=\tfrac12(AB+BA)$.
	\end{enumerate}
\end{definition}

Not every Jordan--Banach algebra admits such an embedding, so the property
deserves its own name. Jordan algebras for which the embedding exists are
singled out below, since it is precisely this class that will later be
identified with the selfadjoint part of a genuine $C^*$-algebra of
observables.

\begin{definition}[JC algebra]\label{def:jc-algebra}
	A JB algebra admitting an embedding as in Definition \ref{def:jb-embedding} is called a \emph{JC algebra}. A
	Jordan algebra of observables admitting such an embedding is called
	\emph{special}, and \emph{exceptional} otherwise.
\end{definition}

\begin{remark}
	Exceptional Jordan algebras exist but are classified by an essentially
	unique finite-dimensional pathological building block, too limited to
	accommodate the observable algebra of a realistic physical system. We
	therefore restrict attention throughout to special Jordan algebras of
	observables, JC algebras, which generate an associative
	$C^*$-algebra $\A$.
\end{remark}

The embedding of Definition~\ref{def:jb-embedding} realizes $\mathcal{O}$
inside a real associative algebra, but a genuine $C^*$-algebra is by
convention a complex object. To hop from the real Jordan algebra
$\mathcal{O}$ to its complex counterpart, one first complexifies, doubling
the real dimension while keeping every operation already defined.

\begin{definition}[Complexification]
	The complexification of a JC algebra $\mathcal{O}$ as in Definition \ref{def:jc-algebra} is
	$\mathcal{O}^{\C}=\{A+iB \ | \ A,B\in\mathcal{O}\}$, with product
	extended by distributivity,
	\[
	(A+iB)\circ(C+iD)=(A\circ C-B\circ D)+i(A\circ D+B\circ C),
	\]
	involution $(A+iB)^*=A-iB$, and states extended by linearity,
	$\omega(A+iB)=\omega(A)+i\,\omega(B)$.
\end{definition}

One checks directly that
$\omega\big((A+iB)\circ(A-iB)\big)=\omega(A^2+B^2)\ge0$ and
$\|A+iB\|=\|A-iB\|$. This yields $\|C\|=\|C^*\|$ for all
$C\in\mathcal{O}^{\C}$. The complexified structure already carries
the positivity and symmetry expected of a $C^*$-algebra. The complexified,
embedded algebra $\A$ generated by $\mathcal{O}$ is thus a unital
$C^*$-algebra with identity $\Id$. Its subset of selfadjoint elements,
which coincides with the original Jordan algebra, is precisely
$\mathcal{O}$, and $\A$ is generated by $\mathcal{O}$.

\subsection*{Lie--Jordan algebras and the associative product}

The embedding just constructed makes a second bilinear operation available
on the ambient algebra $\A$ besides $\circ$, the commutator
$[A,B]:=AB-BA$. The commutator alone does not preserve selfadjointness. 
In fact, fixing a constant $\hbar > 0$,  we can rescale the commutator and
physically identify it with the reduced Planck constant. This rescaling preserves selfadjointness and it is the missing concept needed to recover the full associative
product from $\mathcal{O}$.

\begin{proposition}[Selfadjointness of the rescaled commutator]\label{prop:quantum-comm}
	Let $\mathcal{O}$ be a JC algebra as in Definition \ref{def:jc-algebra} and let $A,B\in\mathcal{O}$. Then $[A,B]^*=-[A,B]$, and consequently
	\[
	\{A,B\} := \frac{1}{i\hbar}[A,B],
	\]
	satisfies $\{A,B\}^*=\{A,B\}$, so that $\{A,B\}\in\mathcal{O}$. The
	bracket $\{\cdot,\cdot\}:\mathcal{O}\times\mathcal{O}\to\mathcal{O}$ is
	thus unambiguously defined on selfadjoint elements.
\end{proposition}
\begin{proof}
	Since $A^*=A$, $B^*=B$,
	$[A,B]^*=(AB-BA)^*=B^*A^*-A^*B^*=BA-AB=-[A,B]$. Hence
	$\{A,B\}^*=\big(\tfrac{1}{i\hbar}[A,B]\big)^*=\tfrac{1}{-i\hbar}[A,B]^*
	=\tfrac{1}{-i\hbar}(-[A,B])=\{A,B\}$.
\end{proof}

Equivalently, $[A,B]=i\hbar\{A,B\}$. Note that $\{\cdot,\cdot\}$ plays the role of the
quantum-mechanical Poisson bracket, and reality of $\{\cdot,\cdot\}$ is
exactly what makes the following axioms consistent with the operational
structure already established on $\mathcal{O}$.

\begin{definition}[Real Lie--Jordan algebra]\label{def:lie-jordan}
	A Jordan--Banach algebra $(\mathcal{O},\circ)$, as in
	Definition~\ref{def:jb-algebra}, equipped with a Lie bracket
	$\{\cdot,\cdot\}$ is a \emph{Lie--Jordan algebra} if, there exists
	$c\in\mathbb{R}$ such that
	\[
	(A\circ B)\circ C-A\circ(B\circ C)=\frac{c}{4}\{\{A,C\},B\}, \qquad
	\{A,B\circ C\}=\{A,B\}\circ C+\{A,C\}\circ B,
	\]
	for all $A,B,C\in\mathcal{O}$.
\end{definition}

Note that, on physical grounds, one sets $c=\hbar^2$,
consistently with $\{\cdot,\cdot\}=\tfrac{1}{i\hbar}[\cdot,\cdot]$ above. The two products of $(\mathcal{O},\circ,\{\cdot,\cdot\})$ are not independent data once the embedding of Definition~\ref{def:jb-embedding} is in place. They recombine into the single associative product of the ambient algebra $\A$, in the precise sense made explicit by the following identity.

\begin{proposition}[Recovery of the associative product]
	Let $\mathcal{O}$ be a JC algebra as in Definition \ref{def:jc-algebra}. For all $A,B\in\mathcal{O}$,
	\begin{equation}
		AB=A\circ B+\tfrac{i\hbar}{2}\{A,B\}. \label{eq:jordan-ass-product}
	\end{equation}
\end{proposition}
\begin{proof}
	A direct computation shows that $A\circ B+\tfrac{i\hbar}{2}\{A,B\}
	=\tfrac12(AB+BA)+\tfrac{i\hbar}{2}\cdot\tfrac{1}{i\hbar}(AB-BA)
	=\tfrac12(AB+BA)+\tfrac12(AB-BA)=AB$.
\end{proof}

This identity can be reversed, and this is the precise sense in which a
Lie--Jordan algebra is associated with a $C^*$-algebra. Starting from the
abstract axioms of Definition~\ref{def:lie-jordan} alone, without
presupposing an ambient associative algebra, one can reconstruct such an
algebra by declaring the right-hand side of the identity above to be the
product.

\begin{theorem}[From Lie--Jordan algebras to $C^*$-algebras]
	\label{thm:lie-jordan-to-cstar}
	Let $(\mathcal{O},\circ,\{\cdot,\cdot\})$ be a real Lie--Jordan algebra, as
	in Definition~\ref{def:lie-jordan}, with $c=\hbar^2$. On
	$\mathcal{O}^{\C}:=\C\otimes\mathcal{O}$, with involution
	$(a\otimes A)^*:=\bar a\otimes A$, the $\C$-linear extension of
	$AB:=A\circ B+\tfrac{i\hbar}{2}\{A,B\}$ is an associative product, and the
	completion of $\mathcal{O}^{\C}$ with respect to its natural norm
	is a $C^*$-algebra $\A$ whose selfadjoint part is $\mathcal{O}$.
\end{theorem}

\noindent A complete proof of this theorem is provided by \cite[p.807]{morettiSpectralTheoryQuantum}.

\begin{remark}
	The algebra $\A$ produced this way is the same pair obtained from the
	JC-algebra embedding of Definition~\ref{def:jb-embedding}, and it is
	exactly the algebra $\A$ of the quantum system of
	Definition~\ref{def:quantum-system} in the next subsection.
\end{remark}

Before stating the formal definition of a quantum system, we summarize the algebraic reconstruction developed in this section. The starting point consists of the operational procedures of preparations and tests, which define the states and the observables. The observable set is structured using rescaling and squaring operations. States are introduced as normalized positive linear functionals, representing the expectation values of measurements. By recording the bounds of these expectations, a norm is defined, and the observable set is completed into a real Banach space. The symmetric Jordan product is introduced to combine observables using only sums and squares. The addition of a real Lie bracket, corresponding to the rescaled commutator, turns the space into a Lie--Jordan algebra. Through Theorem \ref{thm:lie-jordan-to-cstar}, these products generate an associative product on the complexified space. This procedure yields a $C^*$-algebra whose selfadjoint elements coincide with the initial observable set, establishing the mathematical requirements for the axiomatic definition.

\subsection*{The $C^*$-algebraic definition of a Quantum System}\label{sec:alg-quant-sys}

The operational construction of Section \ref{sec:intro-alg-quant} leads, for every physically
reasonable observable set, to a unital $C^*$-algebra $\A$ together with
a set of normalized positive linear functionals on it. This is
the same kind of pair that describes a classical system, except that
commutativity of $\A$ is no longer assumed. We adopt this pair as the
axiomatic definition of a quantum system, and this section collects the
additional facts, automatic continuity of states, the spectrum, and
separation by states, needed to use the definition as a working framework
rather than just formal statement.

\begin{definition}[Quantum system]\label{def:quantum-system}
	A quantum system is defined by the following two sets of data.
	\begin{enumerate}
		\item A unital $C^*$-algebra $\A$ of observables, its selfadjoint
		elements being the operationally accessible observables $\mathcal{O}$ of
		Definition~\ref{def:obs-set}.
		\item A set $S$ of \emph{states}, normalized positive linear functionals
		on $\A$ satisfying $\omega(\Id)=1$ and $\omega(A^*A)\ge0$ for all
		$A\in\A$, such that $S$ is \emph{full}. To wit, it separates the
		elements of $\A$, so that $\omega(A)=\omega(B)$ for all $\omega\in S$
		forces $A=B$. Conversely $\A$ separates $S$.
	\end{enumerate}
\end{definition}

A classical system is recovered as the special case in which $\A$ is abelian, as shown in the previous chapter \ref{sec:classical-alg-rec}. Before this pair can serve as a working framework, it must be checked that positivity alone already forces the continuity that would otherwise need to be postulated separately. The following proposition establishes exactly this automatic regularity.

\begin{proposition}[Automatic continuity and positivity]
	Any positive linear functional $\omega: \A \to \C$ on a unital $C^*$-algebra $\A$, as in Definition \ref{def:unital-c-star-algebra}, is
	automatically continuous in the norm topology,
	\[
	|\omega(A)|\le\|A\|\,\omega(\Id), \qquad \forall A\in \A,
	\]
	and $\omega(A)\ge0$ for all $\omega\in S$ if and only if $A=B^*B$ for some
	$B\in \A$. 
\end{proposition}

This is proved in \cite[Proposition 2.3.11]{bratteliOperatorAlgebrasQuantum1981}.
The spectrum of an observable records, algebraically, the physical values it can legitimately yield upon measurement. We recall its formal definition for a general element of a $C^*$-algebra $\A$.

\begin{definition}[Spectrum]
	Let $\A$ be a $C^*$-algebra in the sense of Definition \ref{def:c-star-algebra}. The spectrum of $A\in \A$ is defined as the set
	$\sigma(A)=\{\lambda\in\C \mid \lambda\Id-A \text{ has no two-sided
		inverse in } \A\}$.
\end{definition}

\noindent The proposition below will show that every value in the spectrum of an observable is attained as an expectation value, but this attainment result holds only for a specific class of elements. We therefore isolate this class first, restricting attention to elements that commute with their own adjoint.

\begin{definition}[Normal elements]\label{def:normal-elements}
	An element $A \in \A$ is called \emph{normal} if $AA^*=A^*A$.
\end{definition}

Every physical observable $A \in \mathcal{O}$ is selfadjoint by definition, and therefore inherently normal. As we will see in the next sections, this framework establishes an isomorphism between the $C^*$-subalgebra generated by $A$ and the algebra of continuous functions on its spectrum, $C(\sigma(A))$. Operationally, this structural correspondence dictates that evaluating the observable behaves exactly like sampling points from $\sigma(A)$. Consequently, the elements of the spectrum represent the only possible deterministic outcomes of a measurement. For normal elements, the functional calculus guarantees that every spectral value corresponds to a sharp physical preparation.

\begin{proposition}[Attainment of spectral values]
	Let $\A$ be a $C^*$-algebra in the sense of Definition \ref{def:c-star-algebra}. If $A\in\A$ is normal as in Definition \ref{def:normal-elements} and $\lambda\in\sigma(A)$, there exists a state $\omega\in S$ such that $\omega(A)=\lambda$.
\end{proposition}

\noindent This Proposition is proved in \cite[Theorem 3.3.6]{murphyCalgebrasOperatorTheory1990}.

\begin{remark}
	The attainment proposition mathematically crystallizes the physical interpretation of the spectrum. Operationally, $\sigma(A)$ constitutes the set of values that a measurement of $A$ will yield with zero variance when the system is prepared in a suitable, dispersion-free state. For a generic state $\omega$, the expectation value $\omega(A)$ is not necessarily a sharp reading but rather a statistical average over the spectrum, thereby falling within the convex hull of $\sigma(A)$.
\end{remark}

\begin{corollary}
	The state space $S$ of Definition \ref{def:state-space} separates all normal elements of $\A$, and consequently all of $\A$. Indeed, any generic element $A\in\A$ can be uniquely decomposed into a complex linear combination of two selfadjoint, and thus normal, elements via the Cartesian decomposition $A=\tfrac12(A+A^*)-\tfrac{i}{2}(iA-iA^*)$.
\end{corollary}

\begin{example}[Continuous spectrum]\label{ex:position-algebra}
	Let $\H=L^2([0,1],dx)$ and let $\A=\mathcal{B}(\H)$ be
	the full algebra of bounded operators on $\H$, unital and
	noncommutative. Define the position observable $Q\in\A$ by
	$(Qf)(x):=xf(x)$ for $f\in\H$. Since $x\in[0,1]$, $Q$ is
	bounded and selfadjoint, hence normal as in
	Definition~\ref{def:normal-elements}. For $\lambda\notin[0,1]$, the
	function $x\mapsto(\lambda-x)^{-1}$ is bounded and continuous on
	$[0,1]$, so multiplication by it provides a two-sided inverse of
	$\lambda\Id-Q$ in $\A$. Conversely, for $\lambda\in[0,1]$ no such
	bounded inverse exists, since $(\lambda-x)^{-1}$ is unbounded near
	$x=\lambda$. Hence $\sigma(Q)=[0,1]$ is a continuous spectrum.
\end{example}
\section{An Introduction to the Hilbert Space Formalism}

The Hilbert space formalism of quantum mechanics was developed through the Dirac-Von Neumann axioms. In particular, quantum states are vectors in a Hilbert space, and observables are operators acting on it. Although we will not develop this framework here, we present the key concepts to later understand how algebraic reconstruction links to the Hilbert space formalism. 

The discussion begins with the inner product and the completion of a pre-Hilbert space into a full Hilbert space, which provides the necessary framework for the rest of the study. From there, we introduce both bounded and unbounded operators, carefully distinguishing them by their domains. In the case of bounded operators, we can construct their adjoint using the Riesz representation theorem as the key tool. A particular point of emphasis is the difference between symmetric and Selfadjoint unbounded operators, as this distinction affects their spectral properties and domain intricacies. We conclude by introducing cyclic vectors, a concept included to lay the groundwork for the GNS theorem developed in the later sections.

\subsection*{Inner Product Spaces and Hilbert Spaces}

We start with a complex vector space endowed with an inner product that yields a notion of length, distance, and orthogonality. A vector space with an inner product is called a pre-Hilbert space.If the normed space is complete under its induced metric, then it is a Hilbert space.

\begin{definition}[Hermitian Inner Product]\label{def:hermitian-ip}
	Let $V$ be a complex vector space. A Hermitian inner product on $V$ is a map $(\cdot,\cdot): V\times V\to \C$ satisfying the following properties for all $v,w,w'\in V$ and $\alpha,\beta\in\C$.
	\begin{enumerate}
		\item $(v,v)\geq 0$.
		\item $(v,v)=0$ if and only if $v=0$.
		\item $( v,\alpha w+\beta w')=\alpha( v,w)+\beta( v,w')$.
		\item $( v,w)=\overline{( w,v)}$.
	\end{enumerate}
\end{definition}

\noindent Note that we adopt the physics convention of linearity in the second argument, so that property 4 combined with property 3 makes the inner product conjugate-linear in the first argument.

\begin{remark}
	Some texts use the opposite convention, linear in the first argument and conjugate-linear in the second. The two conventions differ only by conjugation. We keep the physics convention throughout.
\end{remark}

\noindent The inner product induces a norm through its square root and makes the complex vector space normed.

\begin{proposition}[The Inner Product Induces a Norm]\label{prop:induced-norm}
	Let $V$ be a complex vector space equipped with a Hermitian inner product as per Definition~\ref{def:hermitian-ip}. The map $\|\cdot\|: V\to\mathbb{R}$ defined by $\|v\|=\sqrt{( v,v)}$ is a norm on $V$.
\end{proposition}

\noindent A proof of this statement can be found in \cite[Proposition 3.3]{morettiSpectralTheoryQuantum}. 
\begin{example}
	The space $\C^n$ with $( v,w)=\sum_{k=1}^n\bar v_kw_k$ is a finite-dimensional pre-Hilbert space. The space $\ell^2$ of square-summable sequences, with the analogous inner product, is an infinite-dimensional example used throughout this work.
\end{example}

\noindent Before completeness enters the picture, it is useful to isolate the algebraic content of an inner product space on its own. This intermediate structure will serve as the domain on which the completion procedure of Theorem~\ref{thm:completion} acts.

\begin{definition}[Pre-Hilbert Space]\label{def:pre-hilbert}
	A complex vector space equipped with a Hermitian inner product as per Definition~\ref{def:hermitian-ip} is called a pre-Hilbert space. With the induced norm of Proposition~\ref{prop:induced-norm}, it is a normed space.
\end{definition}

\noindent We can topologically complete a pre-Hilbert space and obtain the central structure used in quantum mechanics.

\begin{definition}[Hilbert Space]\label{def:hilbert}
	A Hilbert space $\H$ is a pre-Hilbert space as per Definition~\ref{def:pre-hilbert} that is complete with respect to its induced norm.
\end{definition}

\begin{example}\label{ex:c01-incomplete}
	The space $C([0,1])$ with $( f,h)=\int_0^1\overline{f(x)}h(x)\,dx$ is a pre-Hilbert space that is not complete, since a sequence of continuous functions approximating a jump discontinuity is Cauchy without a continuous limit. Both $\C^n$ and $\ell^2$ are complete, hence genuine Hilbert spaces.
\end{example}

\noindent The preceding example demonstrates that pre-Hilbert spaces need not be complete. However, completeness is a requirement for many constructions in functional analysis and quantum mechanics. In analogy with the completion of the rationals to the reals, every pre-Hilbert space admits a completion obtained by adjoining the limit points of all Cauchy sequences. To formulate this procedure rigorously, we introduce the notions of dense subsets and linear isometries. The former characterizes how the original space is embedded into its extension, while the latter guarantees that the norm and inner-product structures are preserved under this embedding.

\begin{definition}[Dense Subset]\label{def:dense}
	Let $V$ be a normed space. A subset $W\subseteq V$ is dense in $V$ if for every $v\in V$ and for every $\varepsilon>0$ there exists $w\in W$ with $\|v-w\|<\varepsilon$.
\end{definition}

Both notions will enter the statement of the completion theorem below. Density describes how the original space sits inside its completion, while a linear isometry is precisely the kind of map under which such an embedding must be realized, since it preserves the norm and inner-product structure intact.

\begin{definition}[Linear Isometry]\label{def:isometry}
	Let $V,V'$ be normed spaces. A linear map $J: V\to V'$ is a linear isometry if $\|Jv\|_{V'}=\|v\|_V$ for every $v\in V$.
\end{definition}

\noindent We can show that every pre-Hilbert space embeds densely into a Hilbert space, unique up to unitary isomorphism.

\begin{theorem}[Completion of a Pre-Hilbert Space]\label{thm:completion}
	Let $V$ be a pre-Hilbert space as per Definition~\ref{def:pre-hilbert}. There exists a Hilbert space $\H$ as per Definition~\ref{def:hilbert} and an injective linear isometry $J: V\to\H$ as per Definition~\ref{def:isometry} such that $JV$ is dense in $\H$ as per Definition~\ref{def:dense}. The completion is unique up to unitary isomorphism, and
	\begin{equation}
		( Jv,Jw)_{\H}=( v,w)_V,
	\end{equation}
	for all $v,w\in V$.
\end{theorem}

\noindent A complete proof of this statement is found in \cite[Theorem 3.11]{morettiSpectralTheoryQuantum}.

\subsection*{Bounded Operators on Hilbert Spaces}

In finite-dimensional linear algebra, every linear map is continuous. However, on normed spaces, continuity becomes a non-trivial property, the fundamental equivalence between continuity and boundedness makes the latter a convenient algebraic proxy for analytic regularity. This equivalence leads to the operator norm, which turns the space of operators into a metric setting and it allows us to quantify their size. Then, we will specialize to continuous linear functionals, whose collection forms the dual space. At the end, since many operators of interest are initially defined only on a proper dense subspace, we determine whether they can be extended to the whole space. The bounded extension theorem supplies this, guaranteeing a unique continuous extension whenever the operator is bounded on its dense domain. We start with defining boundedness.

\begin{definition}[Operator Norm and Boundedness]\label{def:op-norm-bdd}
	Let $X,Y$ be normed spaces and let $T\in\mathcal{L}(X,Y)$ with $\mathcal{L}(X,Y)$ the space of linear maps between $X$ and $Y$. The operator $T$ is bounded if there exists $K\geq 0$ such that
	\begin{equation}
		\|Tv\|_Y\leq K\|v\|_X,
	\end{equation}
	for every $v\in X$. Thus the operator norm is $\|T\|=\sup_{v\neq 0} \frac{\|Tv\|_Y}{\|v\|_X}$. 
\end{definition}

\noindent The set of bounded linear operators from $X$ to $Y$ is denoted $\mathcal{B}(X,Y)$. While continuity is inherently a local property at a given point, linearity of an operator forces this local behaviour to propagate globally. Furthermore, a particularly convenient metric for quantifying continuity is the operator norm, which leads to the following fundamental equivalence.

\begin{theorem}[Continuity if and only if Boundedness]\label{thm:cont-iff-bdd}
	Let $T: X\to Y$ be linear, with $X,Y$ normed spaces. The following statements are equivalent.
	\begin{enumerate}
		\item $T$ is continuous at $0\in X$.
		\item $T$ is continuous everywhere on $X$.
		\item $T$ is bounded as per Definition~\ref{def:op-norm-bdd}.
	\end{enumerate}
\end{theorem}

\noindent This theorem is proved in \cite[Theorem 3.11]{morettiSpectralTheoryQuantum}.

\begin{example}\label{ex:shift-diagonal}
	On $\ell^2$, the right shift $(Sv)_1=0$, $(Sv)_{k+1}=v_k$ satisfies $\|Sv\|=\|v\|$, so $S$ is bounded with $\|S\|=1$. On the subspace of finitely supported sequences, $(Dv)_k=kv_k$ is linear but unbounded, since $\|Dv\|=n\|v\|$ for $v$ the $n$-th standard basis vector.
\end{example}

\noindent Having characterized bounded operators between arbitrary normed spaces, we now focus on functionals. These objects lie in dual spaces and provide a coordinate-like description of infinite-dimensional spaces. 

\begin{definition}[Algebraic Dual]\label{def:alg-dual}
	Let $X$ be a complex vector space. The algebraic dual of $X$ is $X^{*}=\mathcal{L}(X,\C)$.
\end{definition}

The algebraic dual imposes no regularity on the functionals it contains, and in infinite dimensions most of them fail to be continuous. Restricting attention to the bounded, hence continuous, functionals singles out the subspace that will actually be useful in the Hilbert space setting.

\begin{definition}[Topological Dual]\label{def:top-dual}
	Let $X$ be a normed space. The topological dual of $X$ is $X'=\mathcal{B}(X,\C)\subseteq X^{*}$, the set of bounded linear functionals as per Definition~\ref{def:alg-dual} and Definition~\ref{def:op-norm-bdd}. 
\end{definition}

\noindent Note also that by Theorem~\ref{thm:cont-iff-bdd}, $X'$ is the space of continuous linear functionals on $X$. In addition, we show that a bounded operator defined on a dense subspace extends uniquely to the whole space.

\begin{proposition}[Extension of Bounded Operators]\label{prop:blt}
	Let $X$ be a normed space, let $Y$ be a Banach space as in Definition \ref{def:banach-space}, and let $W\subseteq X$ be a dense subspace as per Definition~\ref{def:dense}. If $T\in\mathcal{B}(W,Y)$ as per Definition~\ref{def:op-norm-bdd}, there exists a unique $\widetilde T\in\mathcal{B}(X,Y)$ such that $\widetilde T|_W=T$ and $\|\widetilde T\|=\|T\|$.
\end{proposition}

A complete proof can be found in \cite[Proposition 2.47]{morettiSpectralTheoryQuantum}. 
Now, we turn to a more strict class of bounded operators that will be used later.

\begin{definition}[Compact operators]
	\label{def:compact-operators}
	An operator $A\in\BH$, as per Definition~\ref{def:op-norm-bdd}, is
	\emph{compact} if it is a norm limit of finite-rank operators. 
\end{definition}

It can be showed that compact operators form a closed two-sided ideal $\mathfrak{B}_0(\H)\subset\BH$. In addition, in finite dimension $\mathfrak{B}_0(\H)=\BH$. A more detailed discussion can be found in \cite[Chapter 4]{morettiSpectralTheoryQuantum}.

\subsection*{Orthogonal Decomposition of an Hilbert space}

Two vectors in a Hilbert space are orthogonal if their inner product vanishes. Every closed subspace splits $\H$ into two orthogonal pieces, which yields the Riesz representation theorem for continuous linear functionals. 

\begin{definition}[Orthogonal Complement]\label{def:orth-complement}
	Let $\H$ be a Hilbert space as per Definition~\ref{def:hilbert} and let $K\subseteq\H$. The orthogonal complement of $K$ is
	\begin{equation}
		K^{\perp}=\{v\in\H\mid( v,w)=0\text{ for every }w\in K\}.
	\end{equation}
\end{definition}

\noindent The orthogonal complement of a closed subspace $K$ splits $\H$ in two components, namely $K$ and $K^{\perp}$.

\begin{theorem}[Orthogonal Decomposition]\label{thm:orth-decomp}
	Let $\H$ be a Hilbert space as per Definition~\ref{def:hilbert} and let $K\subseteq\H$ be a closed linear subspace. Every $v\in\H$ admits a unique decomposition $v=p+q$, with $p\in K$ and $q\in K^{\perp}$ as per Definition~\ref{def:orth-complement}. Consequently $\H=K\oplus K^{\perp}$, and $K^{\perp}$ is nontrivial whenever $K$ is a proper subspace of $\H$.
\end{theorem}

The proof of this theorem is available in \cite[Theorem 3.13]{morettiSpectralTheoryQuantum}. We recall a particular class of operators that project any vector to a particular subspace of $\H$. Among all these operators, we will refer to the orthogonal ones with just the term projector.

\begin{definition}[Projector]
	\label{def:projector}
	Let $\H$ be a Hilbert space as in Definition \ref{def:hilbert}. A projector on $\H$, also called an orthogonal projection, is a bounded operator $P \in \mathcal B(\H)$ satisfying
	\begin{equation}
		P^2 = P, \qquad P^* = P.
	\end{equation}
\end{definition}

\noindent Rank one projectors arise from a single unit vector, rather than from a higher dimensional closed subspace. Hence, we state here the definition that will be used later.

\begin{definition}[One-dimensional projector]
	\label{def:one-dimensional-projector}
	Let $\H$ be a Hilbert space as in Definition \ref{def:hilbert} and let $\Psi \in \H$ be a unit vector, $\|\Psi\|=1$. The one-dimensional projector, or rank-one projector, associated to $\Psi$ is the projector $P_\Psi$, as per Definition~\ref{def:projector}, defined by
	\begin{equation}
		P_\Psi \Phi := (\Psi,\Phi)\,\Psi, \qquad \forall \Phi \in \H.
	\end{equation}
	Its range $P_\Psi\H$ is the one dimensional subspace $\C \Psi$ spanned by $\Psi$.
\end{definition}

\begin{remark}
	The operator $P_\Psi$ of Definition~\ref{def:one-dimensional-projector} is indeed a projector, as per Definition~\ref{def:projector}. It is self-adjoint, since
	\begin{equation}
		\big(\Phi_1,P_\Psi\Phi_2\big) = (\Psi,\Phi_2)\big(\Phi_1,\Psi\big) = \big(\Phi_1,\Psi\big)(\Psi,\Phi_2) = \big(P_\Psi\Phi_1,\Phi_2\big), \qquad \forall \Phi_1,\Phi_2 \in \H,
	\end{equation}
	using conjugate-linearity of the inner product in its first argument together with $\overline{(\Psi,\Phi_1)} = (\Phi_1,\Psi)$. It is idempotent, since $\|\Psi\|^2=(\Psi,\Psi)=1$ gives
	\begin{equation}
		P_\Psi^2\Phi = (\Psi,\Phi)\,P_\Psi\Psi = (\Psi,\Phi)(\Psi,\Psi)\,\Psi = (\Psi,\Phi)\,\Psi = P_\Psi\Phi, \qquad \forall \Phi \in \H.
	\end{equation}
\end{remark}

\noindent Continuing our discussion, a continuous linear functional $F$ has kernel $\operatorname{Ker}F$, a closed hyperplane. Its orthogonal complement produces the vector representing $F$. This is the content of the Riesz Theorem.

\begin{theorem}[Riesz Representation Theorem]\label{thm:riesz}
	Let $\H$ be a Hilbert space as per Definition~\ref{def:hilbert} and let $\H'$ be its topological dual as per Definition~\ref{def:top-dual}. For every $F\in\H'$ there exists a unique vector $v_F\in\H$ such that
	\begin{equation}
		Fw=( v_F,w),
	\end{equation}
	for every $w\in\H$. The map $F\mapsto v_F$ is a conjugate-linear isometry from $\H'$ onto $\H$.
\end{theorem}

\noindent This is proved in \cite[Theorem 3.16]{morettiSpectralTheoryQuantum}. We propose an example showing an application of Riesz Theorem \ref{thm:riesz}.

\begin{example}\label{ex:riesz-l2}
	On $\H=\ell^2$, fix $a\in\ell^2$ and define $Fw=\sum_{k=1}^{\infty}\bar a_kw_k$. Then $F$ is a bounded linear functional with $\|F\|=\|a\|$, and $Fw=( a,w)$, so the Riesz vector representing $F$ is $a$ itself.
\end{example}

\noindent The Riesz theorem yields the adjoint of a bounded operator, defined in the next section.

\subsection*{The Adjoint Operation and Unbounded Operators}
\label{sec:unbounded}

The adjoint of a bounded operator follows from the Riesz representation Theorem \ref{thm:riesz}, applied to the functional $\psi\mapsto(\varphi,T\psi)$ for fixed $\varphi$. Although we will not show the proof, the existence of $T^*$ follows directly from the Riesz representation theorem. In fact, it guarantees that the bounded functional \(\psi\mapsto(\varphi,T\psi)\) is uniquely represented as \(( T^*\varphi,\psi)\) for some vector \(T^*\varphi\). This unique representability is precisely what ensures the existence and well definedness of the adjoint operator \(T^*\).

\begin{proposition}[Adjoint of Bounded Operators]\label{prop:adjoint-bdd}
	Let $\H_1,\H_2$ be Hilbert spaces as per Definition~\ref{def:hilbert} and let $T\in\mathcal{B}(\H_1,\H_2)$ as per Definition~\ref{def:op-norm-bdd}. There exists a unique operator $T^{*}:\H_2\to\H_1$ such that
	\begin{equation}
		(\varphi,T\psi)_{\H_2}=( T^{*}\varphi,\psi)_{\H_1},
	\end{equation}
	for every $\psi\in\H_1$ and $\varphi\in\H_2$. Moreover $T^{*}\in\mathcal{B}(\H_2,\H_1)$ and $\|T^{*}\|=\|T\|$.
\end{proposition}

This statement is proved in \cite[Proposition 3.38]{morettiSpectralTheoryQuantum}. Comparing $T$ with $T^{*}$ singles out several classes of operators, generalizing normal, Hermitian, and unitary matrices.

\begin{definition}[Special Classes of Bounded Operators]\label{def:special-classes}
	Let $\H$ be a Hilbert space as per Definition~\ref{def:hilbert} and let $T\in\mathcal{B}(\H)$ as per Definition~\ref{def:op-norm-bdd}, with $T^{*}$ as per Proposition~\ref{prop:adjoint-bdd}. The operator $T$ is called
	\begin{enumerate}
		\item normal if $TT^{*}=T^{*}T$,
		\item selfadjoint if $T=T^{*}$,
		\item isometric if $\|T\psi\|=\|\psi\|$ for every $\psi\in\H$,
		\item unitary if it is isometric and surjective,
		\item positive, written $T\geq 0$, if $(\psi,T\psi)\geq 0$ for every $\psi\in\H$.
	\end{enumerate}
\end{definition}

\begin{remark}
	Having set Definition \ref{def:special-classes}, observe the following.
	\begin{itemize}
		\item $T$ is isometric when $T^{*}T=I$, with $I$ the identity on $\H$.
		\item $T$ is unitary when $T^{*}T=TT^{*}=I$.
		\item Every positive operator is selfadjoint, and admits a unique positive square root.
	\end{itemize}
	
\end{remark}

\begin{example}\label{ex:shift-not-unitary}
	The shift $S$ of Example~\ref{ex:shift-diagonal} is isometric but not unitary, since no sequence with first entry $1$ lies in its range.
\end{example}

Bounded operators are defined on the whole space. Position and momentum operators, for example, are not since they satisfy no bound $\|T\psi\|\leq K\|\psi\|$ on their domain. We define such operators as unbounded operators which are densely defined. We will see that we can extend the concept of adjoint also to this other class of operators.

\begin{definition}[Unbounded Operator]\label{def:unbounded-op}
	Let $\H$ be a Hilbert space as per Definition~\ref{def:hilbert}. An unbounded operator on $\H$ is a linear map $T: D(T)\to\H$, where $D(T)$ is a subspace of $\H$ on which no bound as per Definition~\ref{def:op-norm-bdd} is assumed. The graph of $T$ is
	\begin{equation}
		G(T)=\{(\psi,T\psi)\in\H\times\H\mid\psi\in D(T)\}.
	\end{equation}
\end{definition}

\noindent The unique existence of the adjoint of an unbounded operator requires $D(T)$ to be dense.

\begin{definition}[Adjoint of Unbounded Operators]\label{def:adjoint-unbdd}
	Let $T: D(T)\to\H$ be an unbounded operator as per Definition~\ref{def:unbounded-op} with $D(T)$ dense in $\H$ as per Definition~\ref{def:dense}. The adjoint $T^{*}$ has domain
	\begin{equation}
		D(T^{*})=\{\varphi\in\H\mid\exists\,\omega_\varphi\in\H\text{ such that }(\varphi,T\psi)=(\omega_\varphi,\psi) \forall \psi\in D(T)\},
	\end{equation}
	and for $\varphi\in D(T^{*})$ the value $T^{*}\varphi$ is set equal to $\omega_\varphi$.
\end{definition}

\noindent Comparing $T$ with $T^{*}$ yields the unbounded analogues of the bounded case presented in Proposition \ref{prop:adjoint-bdd}.

\begin{definition}[Hermitian Operator]\label{def:hermitian-op}
	Let $T: D(T)\to\H$ be an unbounded operator as per Definition~\ref{def:unbounded-op}. The operator $T$ is Hermitian if
	\begin{equation}
		(\varphi,T\psi)=( T\varphi,\psi),
	\end{equation}
	for all $\varphi,\psi\in D(T)$.
\end{definition}

Hermiticity alone says nothing about the size of the domain $D(T)$, which may fail to be dense. Adding this density requirement singles out the class for which the adjoint itself is guaranteed to be well defined, as the following definition makes precise.

\begin{definition}[Symmetric Operator]\label{def:symmetric-op}
	A Hermitian operator $T$ as per Definition~\ref{def:hermitian-op} whose domain $D(T)$ is dense in $\H$ as per Definition~\ref{def:dense} is called symmetric.
\end{definition}

\noindent Symmetry makes $T^{*}$ well defined.

\begin{definition}[Selfadjoint Operator]\label{def:selfadjoint-unbdd}
	A symmetric operator $T$ as per Definition~\ref{def:symmetric-op} is selfadjoint if $D(T)=D(T^{*})$ and $T=T^{*}$ on this common domain, with $T^{*}$ as per Definition~\ref{def:adjoint-unbdd}.
\end{definition}

\noindent A symmetric operator may still be enlarged to a Selfadjoint one.

\begin{definition}[Essentially Selfadjoint Operator]\label{def:ess-selfadjoint}
	A symmetric operator $T$ as per Definition~\ref{def:symmetric-op} is essentially selfadjoint if $D(T^{*})$ is dense in $\H$ as per Definition~\ref{def:dense} and $T^{**}=T^{*}$.
\end{definition}

The distinction between symmetric, selfadjoint, and essentially selfadjoint operators concerned the interplay of $T$ with its own adjoint. A further classification, familiar from the bounded case, focuses on whether $T$ and $T^*$ commute with each other.

\begin{definition}[Normal Operator, Unbounded Case]\label{def:normal-unbdd}
	An unbounded operator $T$ as per Definition~\ref{def:unbounded-op} with dense domain is normal if $T^{*}T=TT^{*}$ on the common domain where both sides are defined.
\end{definition}

\noindent Another consequence of the Riesz representation theorem is the Hellinger--Toeplitz theorem. It shows that any everywhere-defined symmetric operator on a Hilbert space is automatically bounded. This theorem thus justifies why the unbounded operators that arise in quantum mechanics, such as position and momentum, can never be defined on the entire Hilbert space but only on dense, proper subspaces.

\begin{theorem}[Hellinger-Toeplitz Theorem]\label{thm:helltoep}
	Let $\H$ be a Hilbert space as per Definition~\ref{def:hilbert} and let $T:\H\to\H$ be a Hermitian operator as per Definition~\ref{def:hermitian-op} defined on the whole space $\H$. Then $T$ is bounded and selfadjoint as per Definition~\ref{def:op-norm-bdd} and Definition~\ref{def:selfadjoint-unbdd}.
\end{theorem}

A complete proof can be found in \cite[Theorem 5.14]{morettiSpectralTheoryQuantum}. Symmetric operators that are not selfadjoint exist only on proper dense subspaces.

\begin{example}\label{ex:position-operator}
	On $\H=L^2(\mathbb{R})$, $(Q\psi)(x)=x\psi(x)$ is Hermitian on $D(Q)=\{\psi\in L^2(\mathbb{R})\mid x\psi(x)\in L^2(\mathbb{R})\}$, which is dense. Hence $Q$ is symmetric, and in fact Selfadjoint, since $Q^{*}$ has the same domain and action.
\end{example}

\noindent Within this framework, we introduce the notion of a cyclic vectors. They provide a structural bridge between the abstract operator and its spectral realization. A vector \(\psi\in\H\) is defined cyclic for an operator \(T\) if the  linear span of $\{T^nv\mid n\geq 0\}$ is dense in \(\H\). For unbounded operators, the existence of cyclic vectors guarantees that \(T\) is unitarily equivalent to a multiplication operator on an \(L^2\)-space. Thus, although an unbounded operator cannot be defined everywhere, a single cyclic vector captures all its spectral data.

\begin{definition}[Cyclic Vector for an Operator]\label{def:cyclic-single}
	Let $\H$ be a Hilbert space as per Definition~\ref{def:hilbert} and let $T\in\mathcal{B}(\H)$ as per Definition~\ref{def:op-norm-bdd}. A vector $v\in\H$ is cyclic for $T$ if the linear span of $\{T^nv\mid n\geq 0\}$ is dense in $\H$ as per Definition~\ref{def:dense}.
\end{definition}

\noindent A single cyclic vector need not exist for an arbitrary bounded operator, but the notion becomes considerably more flexible once a whole family of operators is allowed to act jointly. 

\begin{definition}[Cyclic Vector for a Family of Operators]\label{def:cyclic-family}
	Let $\H$ be a Hilbert space as per Definition~\ref{def:hilbert} and let $\mathcal{F}\subseteq\mathcal{B}(\H)$ be a collection of bounded operators as per Definition~\ref{def:op-norm-bdd}. A vector $v\in\H$ is cyclic for $\mathcal{F}$ if the linear span of $\{Av\mid A\in\mathcal{F}\}$ is dense in $\H$ as per Definition~\ref{def:dense}. Taking $\mathcal{F}=\{T^n\mid n\geq 0\}$ one recovers Definition~\ref{def:cyclic-single}.
\end{definition}

\begin{example}\label{ex:cyclic-multiplication}
	On $\H=L^2([0,1])$, with $(T\psi)(x)=x\psi(x)$, the constant function $v=1$ is cyclic for $T$ as per Definition~\ref{def:cyclic-single}, since $\{T^nv\mid n\geq0\}$ spans the polynomials, which is dense in $L^2([0,1])$ by the Weierstrass approximation theorem.
\end{example}
% =====================================================================
% INSERT FOR CHAPTER 2, at the end of Section 2.2, after Example 2.2.35.
% Purely mathematical: no reference to later material of the thesis.
% Labels defined here:
%   def:pvm, def:pvm-integral, def:pvm-support, prop:pvm-calculus, thm:spectral,
%   def:strongly-commuting, thm:joint-spectral, thm:snag, prop:pvm-commutation,
%   prop:pvm-covariance, def:schwartz, prop:schwartz-properties
% Labels used from earlier material (match to the real ones):
%   def:hilbert-space (Def. 2.2.6), def:projector (Def. 2.2.17),
%   def:unbounded-operator (Def. 2.2.25), def:selfadjoint-operator (Def. 2.2.29)
% All \cite{TODO-...} must be matched to the real keys and locations.
% =====================================================================

\subsection*{Projection-Valued Measures and the Spectral Theorem}

For a selfadjoint matrix the spectral theorem splits the space into eigenspaces, on each of which the matrix acts as multiplication by an eigenvalue. For a selfadjoint operator on an infinite dimensional Hilbert space eigenvectors may fail to exist, as for the multiplication by the variable $x$ on $L^2([0,1])$, and the decomposition into eigenspaces has to be replaced by a family of projections labelled by Borel sets. Such a family is called a projection-valued measure. We introduce it on $\mathbb{R}^n$ rather than on $\mathbb{R}$, since the case $n>1$ describes several commuting operators at once, and also the representations of the additive group $\mathbb{R}^n$. We then collect the functional calculus attached to a projection-valued measure and the spectral theorems connecting it with selfadjoint operators and with unitary groups. All the results below are standard, and they are recalled here with references.

\begin{definition}[Projection-valued measure]\label{def:pvm}
	Let $\mathcal{H}$ be a Hilbert space as per Definition~\ref{def:hilbert-space} and let $\mathcal{B}(\mathbb{R}^n)$ be the Borel $\sigma$-algebra of $\mathbb{R}^n$. A \emph{projection-valued measure} on $\mathbb{R}^n$ is a map $E:\mathcal{B}(\mathbb{R}^n)\to B(\mathcal{H})$ such that
	\begin{enumerate}
		\item $E(\Delta)$ is a projector, as per Definition~\ref{def:projector}, for every $\Delta\in\mathcal{B}(\mathbb{R}^n)$,
		\item $E(\mathbb{R}^n)=\mathbb{1}$,
		\item $E\big(\bigcup_{k}\Delta_k\big)\psi=\sum_k E(\Delta_k)\psi$ for every $\psi\in\mathcal{H}$ and every sequence of pairwise disjoint Borel sets $\Delta_k$.
	\end{enumerate}
\end{definition}

\begin{remark}\label{rem:pvm-properties}
	It follows that $E(\emptyset)=0$ and $E(\Delta\cap\Delta')=E(\Delta)E(\Delta')$ for all Borel sets, so that the projections $E(\Delta)$ commute pairwise \cite[TODO-location]{TODO-RS03}. For $\phi,\psi\in\mathcal{H}$ the set function $\mu_{\phi,\psi}(\Delta):=(\phi,E(\Delta)\psi)$ is a complex Borel measure on $\mathbb{R}^n$, and $\mu_\psi:=\mu_{\psi,\psi}$ is a positive measure of total mass $\|\psi\|^2$.
\end{remark}

The measures $\mu_{\phi,\psi}$ allow one to integrate a function against $E$, in the same way as a function is integrated against a scalar measure. The outcome is an operator, which is unbounded in general, and which is the object that the spectral theorem will identify with a function of a selfadjoint operator.

\begin{definition}[Integral against a projection-valued measure]\label{def:pvm-integral}
	Let $E$ be a projection-valued measure on $\mathbb{R}^n$ as per Definition~\ref{def:pvm} and let $f:\mathbb{R}^n\to\mathbb{C}$ be a Borel function. The operator $f(E):=\int f\,dE$ has domain $D(f(E)):=\big\{\psi\in\mathcal{H}\mid\int|f|^2\,d\mu_\psi<\infty\big\}$ and is the unique operator such that
	\begin{equation*}
		(\phi,f(E)\psi)=\int f\,d\mu_{\phi,\psi}, \qquad \phi\in\mathcal{H},\ \psi\in D(f(E)).
	\end{equation*}
\end{definition}

A projection-valued measure may charge only a small part of $\mathbb{R}^n$, and this part controls the spectrum of all the operators $f(E)$. We isolate it before listing the properties of the calculus.

\begin{definition}[Support of a projection-valued measure]\label{def:pvm-support}
	The \emph{support} $\operatorname{supp}E$ of a projection-valued measure $E$ on $\mathbb{R}^n$ is the complement of the union of all open sets $U\subseteq\mathbb{R}^n$ such that $E(U)=0$.
\end{definition}

\begin{remark}
	Since $\mathbb{R}^n$ is second countable, the union of all $E$-null open sets is covered by countably many of them, hence it is itself $E$-null by $\sigma$-additivity. Therefore $\operatorname{supp}E$ is the smallest closed set whose complement is $E$-null, and $E(U)\neq0$ for every open set $U$ meeting $\operatorname{supp}E$.
\end{remark}

We can now state the properties of the functional calculus that are used in the sequel. They are the exact analogue of the rules for functions of a selfadjoint matrix, with the domain issues that unbounded functions introduce.

\begin{proposition}[Functional calculus]\label{prop:pvm-calculus}
	Let $E$ be a projection-valued measure on $\mathbb{R}^n$ and let $f,h:\mathbb{R}^n\to\mathbb{C}$ be Borel functions. Then
	\begin{enumerate}
		\item $f(E)$ is densely defined, closed and normal, with $f(E)^*=\overline{f}(E)$. It is selfadjoint if and only if $f$ is real $E$-almost everywhere, and it is bounded if and only if $f$ is $E$-essentially bounded.
		\item $f(E)h(E)\subseteq(fh)(E)$ with $D(f(E)h(E))=D(h(E))\cap D((fh)(E))$, and $(f+h)(E)$ is the closure of $f(E)+h(E)$, defined on $D(f(E))\cap D(h(E))$.
		\item If $\phi:\mathbb{R}^n\to\mathbb{R}^k$ is Borel, then $\phi_*E(\Delta'):=E(\phi^{-1}(\Delta'))$ is a projection-valued measure on $\mathbb{R}^k$, and $\int f\,d(\phi_*E)=(f\circ\phi)(E)$ for every Borel function $f$ on $\mathbb{R}^k$.
		\item If $W$ is a unitary operator, then $WEW^{-1}$ is a projection-valued measure, and $Wf(E)W^{-1}=f(WEW^{-1})$, with domain $WD(f(E))$.
		\item If $f$ is real, the spectrum of the selfadjoint operator $f(E)$, defined as the set of $\lambda\in\mathbb{C}$ such that $f(E)-\lambda\mathbb{1}$ has no bounded everywhere defined inverse, is the essential range of $f$ with respect to $E$. If $f$ is continuous it equals $\overline{f(\operatorname{supp}E)}$.
	\end{enumerate}
\end{proposition}

\noindent Items 1 and 2 are proved in \cite[TODO-location]{TODO-RS03}, item 3 is the change of variables formula for the measures $\mu_{\phi,\psi}$, item 4 follows from the definition, and item 5 is proved in \cite[TODO-location]{TODO-Mor18}.

The next theorem is the reason why projection-valued measures are relevant. It asserts that every selfadjoint operator is of the form $\int\lambda\,dE(\lambda)$ for a unique measure, so that the calculus above covers all functions of a given operator.

\begin{theorem}[Spectral theorem]\label{thm:spectral}
	Let $A$ be a selfadjoint operator on $\mathcal{H}$, as per Definition~\ref{def:selfadjoint-operator}. There is a unique projection-valued measure $E_A$ on $\mathbb{R}$ such that $A=\int\lambda\,dE_A(\lambda)$, with equality of domains. Moreover $\sigma(A)=\operatorname{supp}E_A$.
\end{theorem}

\noindent This is proved in \cite[TODO-location]{TODO-RS03} and \cite[TODO-location]{TODO-Mor18}, and $E_A(\Delta)$ is called the spectral projection of $A$ on $\Delta$.

To treat several operators at once one needs to say when they can be diagonalized simultaneously. The appropriate notion is the commutation of the spectral projections, which for unbounded operators is stronger than the formal commutation $AB=BA$ on a common domain.

\begin{definition}[Strongly commuting operators]\label{def:strongly-commuting}
	Selfadjoint operators $A_1,\dots,A_n$ on $\mathcal{H}$ are said to \emph{commute strongly} if their spectral projections commute, \textit{i.e.}, if $E_{A_j}(\Delta)E_{A_k}(\Delta')=E_{A_k}(\Delta')E_{A_j}(\Delta)$ for all $j,k$ and all Borel sets $\Delta,\Delta'\subseteq\mathbb{R}$.
\end{definition}

Strong commutation is exactly what allows the one-variable spectral theorem to be upgraded to several variables. The result also tells how sums of commuting operators behave, which is the main use of the next theorem.

\begin{theorem}[Joint spectral theorem]\label{thm:joint-spectral}
	Let $A_1,\dots,A_n$ be selfadjoint operators on $\mathcal{H}$ which commute strongly. There is a unique projection-valued measure $E$ on $\mathbb{R}^n$ such that $A_j=\int\lambda_j\,dE(\boldsymbol{\lambda})$ for all $j$. Conversely, for every projection-valued measure $E$ on $\mathbb{R}^n$ the operators $\int\lambda_j\,dE(\boldsymbol{\lambda})$ are selfadjoint and commute strongly. In this case $A_1+A_2$ is essentially selfadjoint on $D(A_1)\cap D(A_2)$, and its closure is $\int(\lambda_1+\lambda_2)\,dE(\boldsymbol{\lambda})$.
\end{theorem}

\noindent See \cite[TODO-location]{TODO-RS03} and \cite[TODO-location]{TODO-Mor18}.

The same measures describe the strongly continuous unitary representations of the additive group $\mathbb{R}^n$, by a theorem which generalizes the correspondence between one-parameter unitary groups and selfadjoint operators. This version is the one that will be applied to groups of translations.

\begin{theorem}[Stone--Naimark--Ambrose--Godement]\label{thm:snag}
	Let $U:\mathbb{R}^n\to U(\mathcal{H})$ be a strongly continuous unitary representation of the additive group $\mathbb{R}^n$. There is a unique projection-valued measure $E$ on $\mathbb{R}^n$ such that
	\begin{equation*}
		U(\mathbf{t})=\int e^{i\mathbf{t}\cdot\boldsymbol{\lambda}}\,dE(\boldsymbol{\lambda}), \qquad \mathbf{t}\in\mathbb{R}^n.
	\end{equation*}
	Conversely, every projection-valued measure on $\mathbb{R}^n$ defines in this way a strongly continuous unitary representation. The operators $A_j:=\int\lambda_j\,dE(\boldsymbol{\lambda})$ are the selfadjoint generators of the one-parameter groups $t\mapsto U(te_j)$, they commute strongly, and $U(\mathbf{t})=\exp\big(i\sum_jt_jA_j\big)$ in the sense of the calculus of Definition~\ref{def:pvm-integral}.
\end{theorem}

\noindent See \cite[TODO-location]{TODO-SNAG}. For $n=1$ this is the correspondence between strongly continuous one-parameter unitary groups and selfadjoint operators.

Two consequences of the uniqueness in Theorem~\ref{thm:snag} are used repeatedly. The first says that an operator commuting with a unitary representation of $\mathbb{R}^n$ commutes with its spectral measure, and the second describes how a unitary operator conjugating the representation into a twisted version of itself acts on the spectral measure. We prove both, since they reduce to the injectivity of the Fourier transform of complex measures.

\begin{proposition}[Commutation with the spectral measure]\label{prop:pvm-commutation}
	Let $U$ and $E$ be as in Theorem~\ref{thm:snag}. A bounded operator $X$ on $\mathcal{H}$ satisfies $XU(\mathbf{t})=U(\mathbf{t})X$ for all $\mathbf{t}\in\mathbb{R}^n$ if and only if $XE(\Delta)=E(\Delta)X$ for all Borel sets $\Delta\subseteq\mathbb{R}^n$.
\end{proposition}

\begin{proof}
	If $X$ commutes with all $E(\Delta)$, it commutes with the integrals of simple functions, hence with their norm limits, which include every $U(\mathbf{t})$. Conversely, suppose $X$ commutes with all $U(\mathbf{t})$. For $\phi,\psi\in\mathcal{H}$,
	\begin{equation*}
		\int e^{i\mathbf{t}\cdot\boldsymbol{\lambda}}\,d\mu_{X^*\phi,\psi}(\boldsymbol{\lambda})=(X^*\phi,U(\mathbf{t})\psi)=(\phi,U(\mathbf{t})X\psi)=\int e^{i\mathbf{t}\cdot\boldsymbol{\lambda}}\,d\mu_{\phi,X\psi}(\boldsymbol{\lambda}).
	\end{equation*}
	The Fourier transform is injective on finite complex Borel measures on $\mathbb{R}^n$ \cite[TODO-location]{TODO-FourierUniqueness}, hence $\mu_{X^*\phi,\psi}=\mu_{\phi,X\psi}$. Evaluating on $\Delta$ yields $(\phi,XE(\Delta)\psi)=(\phi,E(\Delta)X\psi)$ for all $\phi,\psi$.
\end{proof}

\begin{proposition}[Covariance of the spectral measure]\label{prop:pvm-covariance}
	Let $U$ and $E$ be as in Theorem~\ref{thm:snag}. Let $W$ be a unitary operator on $\mathcal{H}$, let $L$ be an invertible real $n\times n$ matrix and let $\mathbf{c}\in\mathbb{R}^n$, and suppose that
	\begin{equation*}
		W\,U(\mathbf{t})\,W^{-1}=e^{i\mathbf{c}\cdot\mathbf{t}}\,U(L\mathbf{t}), \qquad \mathbf{t}\in\mathbb{R}^n.
	\end{equation*}
	Set $\phi(\boldsymbol{\lambda}):=L^{T}\boldsymbol{\lambda}+\mathbf{c}$. Then $WE(\Delta)W^{-1}=E\big(\phi^{-1}(\Delta)\big)$ for all Borel sets $\Delta$, and $Wf(E)W^{-1}=(f\circ\phi)(E)$ for every Borel function $f$, with equal domains.
\end{proposition}

\begin{proof}
	By item 4 of Proposition~\ref{prop:pvm-calculus}, $\widetilde{E}:=WEW^{-1}$ is a projection-valued measure with $WU(\mathbf{t})W^{-1}=\int e^{i\mathbf{t}\cdot\boldsymbol{\lambda}}\,d\widetilde{E}(\boldsymbol{\lambda})$. On the other hand, using $L\mathbf{t}\cdot\boldsymbol{\lambda}=\mathbf{t}\cdot L^T\boldsymbol{\lambda}$ and item 3 of the same proposition,
	\begin{equation*}
		e^{i\mathbf{c}\cdot\mathbf{t}}U(L\mathbf{t})=\int e^{i\mathbf{t}\cdot\phi(\boldsymbol{\lambda})}\,dE(\boldsymbol{\lambda})=\int e^{i\mathbf{t}\cdot\boldsymbol{\lambda}'}\,d(\phi_*E)(\boldsymbol{\lambda}').
	\end{equation*}
	By the uniqueness in Theorem~\ref{thm:snag}, $\widetilde{E}=\phi_*E$, that is $\widetilde{E}(\Delta)=E(\phi^{-1}(\Delta))$. The statement on functions follows from items 3 and 4 of Proposition~\ref{prop:pvm-calculus}.
\end{proof}

% ---------------------------------------------------------------------

\subsection*{The Schwartz Space}

Position, momentum and the operators built from them are unbounded and defined only on dense subspaces of $L^2(\mathbb{R}^n)$. It is convenient to have a single dense subspace on which all of them, together with all their products, are defined, and which is stable under the natural transformations of $\mathbb{R}^n$. The Schwartz space is such a subspace. It consists of the smooth functions which, together with all their derivatives, decay faster than any inverse polynomial. We define it and we list the properties that make it a common core for the operators of multiplication by the coordinates and of differentiation.

\begin{definition}[Schwartz space]\label{def:schwartz}
	For multi-indices $\alpha,\beta\in\mathbb{N}^n$ write $\mathbf{x}^\beta=x_1^{\beta_1}\cdots x_n^{\beta_n}$ and $\partial^\alpha=\partial_1^{\alpha_1}\cdots\partial_n^{\alpha_n}$. The \emph{Schwartz space} $\mathcal{S}(\mathbb{R}^n)$ is the vector space of functions $\psi\in C^\infty(\mathbb{R}^n;\mathbb{C})$ such that
	\begin{equation*}
		\|\psi\|_{\alpha,\beta}:=\sup_{\mathbf{x}\in\mathbb{R}^n}\big|\mathbf{x}^\beta\,\partial^\alpha\psi(\mathbf{x})\big|<\infty \qquad \text{for all } \alpha,\beta\in\mathbb{N}^n,
	\end{equation*}
	endowed with the topology generated by the seminorms $\|\cdot\|_{\alpha,\beta}$. For $N\in\mathbb{N}$ we set $\mathcal{S}(\mathbb{R}^n;\mathbb{C}^N):=\mathcal{S}(\mathbb{R}^n)\otimes\mathbb{C}^N$, \textit{i.e.}, the space of $\mathbb{C}^N$-valued functions with components in $\mathcal{S}(\mathbb{R}^n)$.
\end{definition}

\begin{proposition}[Properties of the Schwartz space]\label{prop:schwartz-properties}
	The following hold.
	\begin{enumerate}
		\item $C_c^\infty(\mathbb{R}^n)\subseteq\mathcal{S}(\mathbb{R}^n)\subseteq L^2(\mathbb{R}^n)$, the inclusion in $L^2(\mathbb{R}^n)$ is continuous and $\mathcal{S}(\mathbb{R}^n)$ is dense in $L^2(\mathbb{R}^n)$. For $\sigma>0$, $\mathbf{p},\mathbf{x}_0\in\mathbb{R}^n$ the function $\mathbf{x}\mapsto e^{i\mathbf{p}\cdot\mathbf{x}}e^{-|\mathbf{x}-\mathbf{x}_0|^2/4\sigma^2}$ belongs to $\mathcal{S}(\mathbb{R}^n)$.
		\item The maps $\psi\mapsto x_j\psi$, $\psi\mapsto\partial_j\psi$, $\psi\mapsto\psi(\cdot-\mathbf{a})$, $\psi\mapsto e^{i\mathbf{b}\cdot\mathbf{x}}\psi$ and $\psi\mapsto\psi\circ R^{-1}$, for $\mathbf{a},\mathbf{b}\in\mathbb{R}^n$ and $R\in O(n)$, are continuous linear maps of $\mathcal{S}(\mathbb{R}^n)$ into itself. The same holds for the multiplication by any smooth function all of whose derivatives are polynomially bounded.
		\item The Fourier transform $(\mathcal{F}\psi)(\mathbf{p})=(2\pi)^{-n/2}\int e^{-i\mathbf{p}\cdot\mathbf{x}}\psi(\mathbf{x})\,d\mathbf{x}$ is a homeomorphism of $\mathcal{S}(\mathbb{R}^n)$ onto itself and it extends to a unitary operator on $L^2(\mathbb{R}^n)$.
		\item If $\psi\in L^2(\mathbb{R}^n)$ and all the weak derivatives satisfy $\mathbf{x}^\beta\partial^\alpha\psi\in L^2(\mathbb{R}^n)$ for all $\alpha,\beta\in\mathbb{N}^n$, then $\psi\in\mathcal{S}(\mathbb{R}^n)$.
		\item The operators $Q_j\psi=x_j\psi$ and $P_j\psi=-i\partial_j\psi$ on their natural domains, the latter in the weak sense, are selfadjoint, and $\mathcal{S}(\mathbb{R}^n)$ is invariant under both and is a core for each of them.
	\end{enumerate}
\end{proposition}

\noindent Items 1 to 3 are proved in \cite[TODO-location]{TODO-RS03}. Item 4 follows from the Sobolev embedding theorem applied to each function $\mathbf{x}^\beta\psi$ \cite[TODO-location]{TODO-Sobolev}, and item 5 is proved in \cite[TODO-location]{TODO-RS03}.
\end{document}